% % % % % % % % % % % % % % 
%
%     Skript zu NUMERIK I
%           WS14/15
%    von Prof. Dr. Blank
% Universität Regensburg
%
%
%	Kap. 3: Fehleranalyse
%
% % % % % % % % % % % % % %
\chapter{Fehleranalyse} \index{Fehler}\label{3}
\marginpar{15.10.2014}
\begin{align*}
  \overset{x+\varepsilon \text{ statt } x}{\framebox[3cm]{Eingabe}} \longrightarrow 
  \overset{\underset{\text{\tiny (z.B. durch Rundung)}}{\widetilde{f} \text{ statt } f}}{\framebox[3cm]{Algorithmus}} \longrightarrow
  \overset{\widetilde{f}(x+\varepsilon) \text{ statt } f(x)}{\framebox[3cm]{Resultat\phantom{g}}}
\end{align*}\\

Bei der Fehleranalyse liegt das Hauptaugenmerk auf
\begin{itemize}
\item[] \textbf{Eingabefehler}\\ z.B.Rundungsfehler, Fehler in Messdaten, Fehler im Modell (falsche Parameter)
\item[] \textbf{Fehler im Algorithmus} \\ z.B. Rundungsfehler durch Rechenoperationen, Approximationen \\
  (z.B. Ableitung durch Differenzenquotient oder die Berechnung von Sinus durch abgebrochene Reihenentwicklung)
  \\
\item[\textit{1. Frage}] Wie wirken sich Eingabefehler auf das Resultat unabhängig vom gewählten Algorithmus aus?
\item[\textit{2. Frage}]Wie wirken sich (Rundungs-)Fehler des Algorithmus aus?\\
  Und wie verstärkt der Algorithmus Eingabefehler?
\end{itemize}


\sectione{Zahlendarstellung und Rundungsfehler} \label{3.1} \index{Fehler} \index{Rundungsfehler}\index{Zahlendarstellung}
Auf (Digital-)Rechnern können nur endlich viele Zahlen realisiert werden. \\
Die wichtigsten Typen sind: 
\begin{itemize}
\item \textbf{ganze Zahlen}  (integer)\index{integer}:
  \begin{align*}
    z&=\pm \sum_{i=0}^{m}z_i\beta_i & \text{mit }
                                      \begin{array}{l@{\,}l}
                                        \beta &= \text{Basis des Zahlensystems (oft $\beta=2$)} \\
                                        z_i &\in \{0, \cdots \beta-1\}
                                      \end{array}
  \end{align*}
\item \textbf{Gleitpunktzahlen} (floating point) \index{floating point}
\end{itemize}

% \subsectione{Definition: Gleitkommazahl} 
\begin{Defe}
  \label{3.1.1} \index{Gleitkommazahl}
  Eine Zahl $x\in\Q$ mit einer Darstellung
  \begin{align*}
    x&=\sigma \cdot(a_1 . a_2 \cdots a_t)_{\beta}\cdot \beta^e 
       = \sigma\beta^e \cdot \sum_{\nu=1}^{t}a_\nu \beta^{-\nu+1}\\\\
     &\quad\begin{array}{ll}
             \beta\in\N & \text{Basis des Zahlensystems}\index{Basis}\\
             \sigma \in\{\pm 1\} &\text{Vorzeichen} \\
             m = (a_1 . a_2 \cdots a_t)_{\beta} &\text{Mantisse}\index{Mantisse}\\
             \phantom{m}= \sum_{\nu=1}^{t}a_\nu \beta^{-\nu+1} \\
             a_i \in\{0,\cdots , \beta-1\}&\text{Ziffern der Mantisse}\\
             t\in\N&\text{Mantissenlänge} \\
             e\in\Z &\text{mit }e_{min}\leq e \leq e_{max} \text{ Exponent}
           \end{array}
  \end{align*}
  heißt \textbf{Gleitkommazahl} mit $t$ Stellen und Exponent $e$ zur Basis $b$. \\
  Ist $a_1\neq 0$, so heißt $x$ \textbf{normalisierte Gleitkommazahl}\index{normalisierte Gleitkommazahl}.
\end{Defe}

% \subsectione{Bemerkung} 
\begin{Beme}
  \label{3.1.2}~
  \begin{enumerate}[a)]
  \item 0 ist keine normalisierte Gleitkommazahl, da $a_1 =  0$ ist.
  \item $a_1\neq 0$ stellt sicher, dass die Gleitkommadarstellung eindeutig ist.
  \item In der Praxis werden auch nicht-normalisierte Darstellungen verwendet.
  \item Heutige Rechner verwenden meist $\beta =2$, aber auch $\beta=8, \beta=16$.
  \end{enumerate}
\end{Beme}

\subsectione{Bit-Darstellung zur Basis 2}
\label{3.1.3}
Bit-Darstellung nach IEEE-Standard 754 von floating point numbers \\
Sei die Basis $\beta=2$.

\begin{tabular}{l@{}cccc@{}}
  & Speicherplatz & $t$ & $e_{min}$ & $e_{max}$ \\
  \cmidrule{2-5}
  einfache Genauigkeit (float) \index{floating point} & 32bits = 4Bytes & 24 &-126 & 127 \\
  doppelte  Genauigkeit (double)~~\index{double} & 64bits = 8Bytes& 52 & -1022 & 1023
\end{tabular}\\

Darstellung im Rechner (Bitmuster) für float:
\begin{gather*}
  \floatbox{s}{b_0\cdots b_7}{a_2\cdots\cdots a_{24}}\\
  \text{(Da $a_1\neq 0$, also $a_1=1$ gilt, wird $a_1$ nicht gespeichert)}
\end{gather*}

Interpretation ($s,b,a_i\in\{0,1\} \forall i$)
\begin{itemize}
\item $s$ Vorzeichenbit: 
  $\quad \sigma=(-1)^s\Rightarrow 
  \begin{array}{l}
    \sigma(0)=1 \\
    \sigma(1)=-1
  \end{array} $
\item $b=\sum_{i=0}^{7}b_i\cdot2^i \in \{1, \cdots, 254\}$ speichert den Exponenten mit \\
  $ \quad e = b-\underbrace{127}_\text{Basiswert}$ (kein Vorzeichen nötig) \\
  Beachte: $b_0=\cdots=b_7=1$ sowie $b_0=\cdots=b_7=0$ sind bis auf Ausnahmen keine gültigen Exponenten
\item $m=(a_1.a_2\cdots a_{24})=1+\sum_{\nu=2}^{24}a_{\nu}2^{1-\nu}$ stellt die Mantisse dar, $a_1=1$ wird nicht abgespeichert.
\item Besondere Zahlen per Konvention:
  \begin{itemize}
  \item[$x=0$:] $s$ bel., $b=0$, $m=1 \quad \floatbox{s}{0\cdots0}{0\cdots0}$
  \item[$x=\pm\infty$:]  $s$ bel., $b=255$, $m=1  \quad \floatbox{s}{1\cdots1}{0\cdots0}$
  \item[$x=$NaN] $s$ bel., $b=255$, $m\neq 1$
  \item[$x=(-1)^s$] $s$ bel., $b=0$, $m\neq 1$ und x hat die Form $x=(0+\sum_{\nu=2}^{24}a_{\nu}\cdot 2^{1-\nu})\cdot 2^{126}$ (\enquote{denormalized} number)
  \end{itemize}
\end{itemize}

%--------------------------------------------------------------------

\marginpar{20.10.2014}
\begin{align*}
  \intertext{Betragsmäßig \textbf{größte Zahl}:}
  \floatbox{0}{01\cdots 1}{ 1\cdots\cdots 1} && 
                                                x_{max} = (2-2^{-23})\cdot 2^{127}  & \approx 3,4 \cdot 10^{38}
                                                                                      \intertext{Betragsmäßig \textbf{kleinste Zahl}:}
                                                                                      \floatbox{0}{0\cdots 0}{ 0\cdots\cdots 01} && 
                                                                                                                                    x_{min} = (2-2^{-23})\cdot 2^{-126} = 2^{-149}  & \approx 1,4 \cdot 10^{-45}
\end{align*}

\subsectione{Verteilung der Maschinenzahlen} \label{3.1.4}
ungleichmäßig im Dezimalsystem, z. B.
\begin{align*}
  x &= \pm a_1 . a_2 a_3 \cdot 2^e  & -2\leq & e\leq 1 & a_i & \in \{0,1\} 
\end{align*}
\begin{figure}
  \parbox{\linewidth}{
    \centering
    \includegraphics[width=2cm]{images/image_missing.jpg}
  }
  \caption{Ungleichmäßige Verteilung der Maschinenzahlen im Dezimalsystem}
\end{figure}
ist im Dualsystem gleichmäßig verteilt.

% \subsection{Bezeichnungen}
\begin{Defe}
  \label{3.1.5}~
  \begin{itemize}
  \item[\textbf{overflow}] es ergibt sich eine Zahl, die betragsmäßig größer ist als die größte maschinendarstellbare Zahl
  \item[\textbf{underflow}] entsprechend, betragsmäßig kleiner als die kleinste positive Zahl
  \end{itemize}
  Bsp.: overflow beim integer $b=e+127$
  \begin{align*}
    \begin{array}{rrr@{}}
      b &= 254                                &11111110 \\
        &+  \phantom{24}3 &00000011 \\
      \cline{3-3} %
      b+3 = 257 \text{ mod } 2^8  &=\phantom{24}1& \xout{1}00000001 
    \end{array}	  
  \end{align*}
\end{Defe}


\subsectione{Rundungsfehler} \label{3.1.6}
Habe $x\in \R $ die normalisierte Darstellung
\begin{align*}
  x &= \sigma \cdot \beta^e (\sum_{\nu=1}^{t}a_{\nu}\beta^{1-\nu} + \sum_{\nu=t+1}^{\infty}a_{\nu}\beta^{1-\nu} ) \\
    &= \sigma \cdot \beta^e (\sum_{\nu=1}^{t}a_{\nu}\beta^{1-\nu} + \beta^{1-t}\sum_{l=1}^{\infty}a_{t+l}\beta^{-l} )
\end{align*}
mit $e_{min} \leq e \leq e_{max}$, dann wird mit $fl(x)$ die gerundete Zahl bezeichnet, wobei $fl(x)$ 
eindeutig gegeben ist durch die Schranke an den \textbf{absoluten Rundungsfehler} \index{Fehler!absoluter Rundungsfehler}
\begin{align*}
  | fl(x) - x | \leq \begin{cases}
    \frac{1}{2}\beta^{e+1+t} & \text{bei symmetrischem Runden}\\
    \beta^{e+1+t}                    & \text{bei Abschneiden}
  \end{cases} \quad .
\end{align*}
Für die \textbf{relative Rechengenauigkeit} \index{Genauigkeit!relative Rechengenauigkeit}folgt somit 
\begin{align*}
  \frac{| fl(x) - x | }{|x|} & \leq \begin{cases}
    \frac{1}{2}\beta^{1-t} & \text{bei symmetrischem Runden}\\
    \beta^{1-t}                    & \text{bei Abschneiden}
  \end{cases} \quad .
\end{align*}
Die \textbf{Maschinengenauigkeit} \index{Genauigkeit!Maschinengenauigkeit} des Rechners ist daher durch 
\begin{align*}
  eps &= \beta^{1-t} & \text{(für float}\approx 10^{-7}  \text{, für double} \approx10^{-16} )
\end{align*}
gegeben.

Die Mantissenlänge bestimmt also die Maschinengenauigkeit. Bei einfacher Genauigkeit ist $fl(x)$ bis auf ungefähr 7 signifikante Stellen genau. \\
Im Folgenden betrachten wir symmetrisches Runden und definieren daher
\[ \tau \coloneqq \frac{1}{2}eps\]
Weiterhin gilt:
\begin{enumerate}[a)]
\item Die kleinste Zahl am Rechner, welche größer als 1 ist, ist
  \[ 1 + eps \]
\item Eine Maschinenzahl x repräsentiert eine Eingabemenge
  \[  E(x) = \{\widetilde{x} \in \R : |\widetilde{x}-x| \leq \tau|x|\} \] \\
  \begin{figure}
    \parbox{\linewidth}{
      \centering
      \includegraphics[width=2cm]{images/image_missing.jpg}
    }
    \caption{Eingabemenge einer Maschinenzahl}
  \end{figure}
\end{enumerate}

% \subsection{Bemerkung} 
\begin{Beme}
  \label{3.1.7}
  Gesetze der arithmetischen Operationen gelten i.A. nicht, z.B.
  \begin{itemize}
  \item 	$x$ Maschinenzahl $\quad \Rightarrow fl(x+\nu) = x \text{     für }|\nu| < \tau |x|$
  \item Assoziativ- und Distributivgesetze gelten nicht, z.B. für $\beta = 10, \, t=3, \, a=0,1 ,\, b= 105 , \, c= -104$ gilt:
    \begin{align*}
      fl(a+fl(b+c)) &= 1,1 \\
      fl(fl(a+b)+c) &= fl(fl(105,1) + (-104) ) \\
                    &= fl(105-104) \\
                    &= 1 \quad \lightning
    \end{align*}
  \item[ $\Rightarrow$] Für einen Algorithmus ist die Reihenfolge der Operationen wesentlich!
    Mathematisch äquivalente Formulierungen können zu verschiedenen Ergebnissen führen.
  \end{itemize}
\end{Beme}

\subsectione{Auslöschung von signifikanten Stellen} \label{3.1.8}
Sei $x=9,995\cdot 10^{-1}, y=9,984 \cdot 10^{-1}$. Runde auf drei signifikante Stellen und berechne $x-y$:
\begin{align*}
  \widetilde{f}(x,y) &\coloneqq fl(fl(x)- fl(y)) = fl(1,00\cdot 10^0 - 9,98\cdot 10^{-1}) \\
                     &= 	fl(0,02\cdot 10^{-1}) \\
                     &= fl(2,00 \cdot 10^{-3}) \\
  f(x,y)  &\coloneqq x-y \\
                     &\coloneqq 0,0011 = 1,1\cdot 10^{-3}
                       \intertext{Daraus ergibt sich der relative Fehler}
                       \frac{|\widetilde{f}(x,y)-f(x,y)|}{|f(x,y)|}
                     &= \frac{|2\cdot 10^{-3}- 1,1\cdot 10^{-3}|}{|1,1\cdot 10^{-3}}
                       = 82\%
\end{align*}
Der Grund hierfür ist, dass das Problem der Substraktion zweier annähernd gleich großer Zahlen
schlecht konditioniert ist.\\

\textbf{Zwei Regeln:}
\begin{enumerate}[1)]
\item Umgehbare Substraktion annähernd gleich großer Zahlen vermeiden!
\item Unumgängliche Substraktion möglichst an den Anfang des Algorithmus stellen! (siehe später)
\end{enumerate}

% 2.2
% -----------------------------------------------------------------------------------------------------------------------------------------
\sectione{Kondition eines Problems} \label{3.2}
Es wird das Verhältnis 
\begin{gather*}
  \frac{\text{Ausgabefehler}}{\text{Eingabefehler}}
\end{gather*}
untersucht.

% \subsectione{Definition: Problem} 
\begin{Defe}
  \label{3.2.1} \index{Problem}
  Sei $f: U \subseteq \R^n \rightarrow \R^m$ mit $U$ offen und sei $x\in U$.
  Dann bezeichne $(f, x)$ das Problem, zu einem gegebenen $x$ die Lösung $f(x)$ zu finden.
\end{Defe}

% \subsectione{Definition: absoluter und relativer Fehler} 
\begin{Defe}
  \label{3.2.2} \index{Fehler}
  Sei $x\in\R^n$ und $\widetilde{x} \in \R^n$ eine Näherung an $x$. Weiterhin sei $\|\cdot\|$ eine Norm auf $\R^n$.
  \begin{itemize}
  \item[a)] $\nn{\widetilde{x} - x}$ heißt \textbf{absoluter Fehler} \index{Fehler!absoluter}
  \item[b)] $\frac{\nn{\widetilde{x} - x}}{\nn{x}}$ heißt \textbf{relativer Fehler}\index{Fehler!relativer}
  \end{itemize}
  Da der relative Fehler skalierungsinvariant ist, d.h. unabhänging von der  Wahl von $x$ ist, ist dieser i.d.R. von größerem Interesse.
  Beide Fehler hängen von der Wahl der Norm ab!
  Häufig werden Fehler auch komponentenweise gemessen:
  \begin{align*}
    \text{Für } i=1,\cdots , n : && |\widetilde{x}_i - x_i | & \leq \delta & \text{ (absolut)} \\
                                 && |\widetilde{x}_i - x_i | &\leq \delta |x_i| & \text{ (relativ)}
  \end{align*}
\end{Defe}


% \subsectione{Wiederholung: Normen} 
\begin{Wdhe}[Normen]
  \label{3.2.3}\index{Norm}
  \begin{figure}
    \centering
    \begin{tikzpicture}[scale=2, line cap=round, x=1.0cm,y=1.0cm]
      \clip (-1.5,-1.5) rectangle (2.3,1.5);
      \draw[->] (-1.3,0) -- (1.3,0) node [anchor=north west]{$x$};
      \draw[->] (0,-1.3) -- (0,1.3) node [anchor=south west]{$y$};
      \draw [dashed] (-1,-1) rectangle (1,1);
      \draw (0,0) circle (1cm);
      \draw [dotted] (0,1) -- (1,0) -- (0,-1) -- (-1,0) -- cycle;
      \draw [dotted] (1.3,1.3) -- (1.5,1.3) node [anchor=west]{$\|x\|_{1} = 1$};
      \draw (1.3,1.0) -- (1.5,1.0) node [anchor=west]{$\|x\|_{2} = 1$};
      \draw [dashed] (1.3,0.7) -- (1.5,0.7) node [anchor=west]{$\|x\|_{\infty} = 1$};
    \end{tikzpicture}
    \caption{Sphären mit gleichem Normbetrag}
  \end{figure}
  
  \begin{align*}
    \text{Summennorm ($l_1$-Norm):} &&	\nn{x}_1 &\coloneqq \sum_{i=1}^{n}|x_i| 
                                                   \index{Norm!Summennorm}\\
    \\
    \text{Euklidische Norm ($l_2$-Norm):} &&	\nn{x}_2 &\coloneqq \sqrt{\sum_{i=1}^{n}|x_i|^2}
                                                           \index{Norm!Euklidische Norm} \\
    \\
    \text{Maximumsnorm ($l_\infty$-Norm):} &&\nn{x}_\infty &\coloneqq \max\{|x_i| : i=1, \cdots n\} \\
    \index{Norm!Maximumsnorm}
    \\
    \text{Hölder-Norm ($l_p$-Norm):} &&	\nn{x}_p &\coloneqq 
                                                   \left(\sum_{i=1}^{n}|x_i|^p\right)^{\frac{1}{p}} 
                                                   \index{Norm!Hölder-Norm}
  \end{align*}
\end{Wdhe}



% \subsectione{Definition: Matrixnorm} 
\begin{Defe}
  \label{3.2.4}
  Auf dem $\R^n$  sei die Norm $\nn{\,\cdot\,}_a$ und auf dem $\R^m$ die Norm $\nn{\,\cdot\,}_b$ gegeben.
  Dann ist die zugehörige \textbf{Matrixnorm} \index{Norm!Matrixnorm} gegeben durch:
  \begin{align}
    \nn{A}_{a,b} &\coloneqq \sup_{x\neq 0} \frac{\nn{Ax}_b}{\nn{x}_a} \\ \nonumber
                 &= \sup_{\nn{x}_a=1} \nn{Ax}_b \label{III.2.1} 
  \end{align}
  Also ist   $\nn{A}_{a,b}$ die kleinste Zahl $c>0$ mit
  \begin{gather*}
    \nn{Ax}_b  \leq c\nn{x}_a \quad\quad \forall x\in \R^n
  \end{gather*}
\end{Defe}

% \subsectione{Definition: Frobeniusnorm, p-Norm, Verträglichkeit} 
\begin{Defe}
  \label{3.2.5}
  Sei $A\in \R^{m\times n}$.
  \begin{enumerate}[a)]
  \item \textbf{Frobeniusnorm} (Schurnorm):
    $ \quad \nn{A}_F \coloneqq \sqrt{\sum_{i=1}^{m}\sum_{j=1}^{n}|a_{ij}^2|}
    \index{Norm!Frobeniusnorm}$
  \item \textbf{p-Norm}: 
    $\quad \nn{A}_p \coloneqq \nn{A}_{p,p}
    \index{p-Norm}$
  \item Eine Matrixnorm heißt \textbf{verträglich} \index{Norm!verträglich} mit den Vektornormen 
    $\nn{\,\cdot\,}_a, \nn{\,\cdot\,}_b$, falls gilt
    \footnote{ Beachte: $\nn{A}_{a,b}$ ist die kleinste Norm im Gegensatz zu $\nn{A}$, welche hier beliebig ist.}:
    \begin{gather*}
      \nn{Ax}_b \leq \nn{A}\cdot \nn{x}_a \quad \forall x\in \R^n
    \end{gather*}
  \end{enumerate}
\end{Defe}

% \subsectione{Bemerkungen} 
\begin{Beme}~
  \label{3.2.6}
  \begin{enumerate}[a)]
  \item Die Normen $\nn{\,\cdot\,}_F$ und $\nn{\,\cdot\,}_p$ sind \textbf{submultiplikativ} \index{Norm!submultiplikative}, d.h.
    \begin{gather*}
      \nn{A\cdot B} \leq \nn{A}\cdot\nn{B}
    \end{gather*}
  \item Die Norm $\nn{\,\cdot\,}_{1,1}$ wird auch \textbf{Spaltensummennorm}\index{Norm!Spaltensummennorm} genannt:
    \begin{gather*}
      \nn{A}_1 = \max_{1\leq j \leq n}\sum_{i=1}^{m}|a_{ij}|
    \end{gather*}
    Sie ist das Maximum der Spaltensummen\footnote{Beweis: siehe Übungsblatt 3}.
  \item Die Norm $\nn{\,\cdot\,}_{\infty, \infty}$ wird auch \textbf{Zeilensummennorm} \index{Zeilensummennorm}
    genannt\footnote{Beweis: siehe Übungsblatt 3}:
    % not sure, why \footref won't work here ...
    \begin{gather*}
      \nn{A}_\infty = \max_{1\leq i \leq m}\sum_{j=1}^{n}|a_{ij}|
    \end{gather*}
  \item Die Frobeniusnorm $\nn{\,\cdot\,}_F$ ist verträglich mit der euklidischen Norm $\nn{\,\cdot\,}_2$
  \item Die Wurzeln aus den Eigenwerten von $A^TA$ heißen \textbf{Singulärwerte $\sigma_i$} \index{Singulärwert} von A.
    Mit ihnen kann die $\nn{\,\cdot\,}_{2,2}$ Norm dargestellt werden\footnote{Beweis: siehe Übungsblatt 3}:
    \begin{align*}
      \nn{A}_2 &= max \{\sqrt{\mu} : A^TA\cdot x = \mu x\text{ für ein }x\neq 0 \} \\
               & = \sigma_{max}
    \end{align*}
  \end{enumerate}
\end{Beme}


% ultimate evil hack to go along with numeration
% \minisec{\Large3.2 a) Normweise Konditionsanalyse} \label{3.2a}\vspace{1eM}
\extrasection{a)}{Normweise Konditionsanalyse}
%% Alternative to add it to table of contents:
% \renewcommand{\thesubsection}{\thesection.a)}
% \setkomafont{subsection}{\Large}
% \subsectione{NORMWEISE KONDITIONSANALYSE}
% \renewcommand{\thesubsection}{\thesection.\arabic{subsection}}
% \addtocounter{subsection}{-1}
% \setkomafont{subsection}{\large}

% \subsectione{Definition: absolute und relative normweise Kondition}

%---------------------------------------------

\begin{Defe}
	\marginpar{22.10.2014}
  \index{normweise Kondition}
  Sei $(f,x)$ ein Problem mit $f:U\subset \R^n \rightarrow \R^m$
  und $\nn{\,\cdot\,}_a$ auf $\R^n$ und $\nn{\,\cdot\,}_b$ auf $\R^m$ eine Norm.
  \begin{enumerate}[a)]
  \item Die \textbf{absolute normweise Kondition}\index{Kondition!normweise, absolut} eines Problems $(f,x)$ ist die kleinste Zahl 
    $\kappa _{abs} > 0 $ mit
    \begin{align}
      \nn{f(\widetilde{x})-f(x)}_b &\leq \kappa _{abs}(f,x) \nn{\widetilde{x}-x}_a
                                     + o\left(\nn{\widetilde{x}-x}_a\right) \label{III.2.2} \\\nonumber
      \Bigl(f(\widetilde{x})- f(x) 
                                   &=\underbrace{ f'(x)(\widetilde{x}-x)\pm o\left(\nn{\widetilde{x}-x}\right)}_{Taylorentwicklung}
                                     \quad \text{für }\widetilde{x}\rightarrow x 
                                     \Bigr)
    \end{align}
  \item Die \textbf{relative normweise Kondition}\index{Kondition!normweise. relativ} eines Problems $(f,x)$  mit $x\neq 0, f(x) \neq 0$
    ist die kleinste Zahl 
    $\kappa _{rel} > 0 $ mit
    \begin{align}
      \frac{	\nn{f(\widetilde{x})-f(x)}_b }{\nn{f(x)}_b}
      &\leq \kappa _{rel}(f,x)\frac{ \nn{\widetilde{x}-x}_a}{\nn{x}_a}
        + 
        o\left(\frac{\nn{\widetilde{x}-x}_a}{\nn{x}_a}\right) \label{III.2.3}
      &&	\text{für } \widetilde{x} \rightarrow x
    \end{align}
  \item Sprechweise:
    \begin{itemize}\index{Kondition!gut/schlecht konditioniert}
    \item falls $\kappa$ \enquote{klein} ist, ist das Problem \enquote{gut konditioniert}
    \item falls $\kappa$ \enquote{groß} ist, ist das Problem \enquote{schlecht konditioniert}
    \end{itemize}
  \end{enumerate}
\end{Defe}

% \subsectione{Lemma} 
\begin{Leme}\label{3.2.8}
  Falls $f$ differenzierbar ist, gilt
  \begin{gather}
    \kappa_{abs}(f,x) = \nn{Df(x)}_{a,b} \label{III.2.4}
  \end{gather}
  und für $f(x) \neq 0$
  \begin{gather}
    \kappa_{rel}(f,x) = \frac{\nn{x}_a}{\nn{f(x)}_b}\cdot \|Df(x)\|_{a,b} \label{III.2.5}
  \end{gather}
  wobei $Df(x)$ die Jakobi-Matrix bezeichnet.
\end{Leme}

% \subsectione{Beispiel: Kondition der Addition}
\begin{Bspe}[Kondition der Addition]
  \label{3.2.9} \index{Kondition!Addition}
  $f(x_1, x_2) \coloneqq x_1 +x_2 , \, f:\R^2 \rightarrow \R$. \\
  Wähle $l_1$-Norm auf $\R^2$ (und $\R$)
  \begin{align*}
    Df(x_1, x_2) \, =(\nabla f^T) \, &= (\frac{\partial}{\partial x_1}f, \frac{\partial}{\partial x_2}f )\\
                                     &= (1,1) && \text{(Matrix!)}
  \end{align*}
  damit
  \begin{align*}
    \kappa_{abs} (f,x)&= \nn{Df(x)}_{1,1} && \text{(Matrix-Norm!!)}\\
                      &= \nn{Df(x)}_1 \\
                      &=1 \\
    \kappa_{rel} (f,x) &= \frac{\nn{x}_1}{\nn{f(x)}_1} \cdot \nn{Df(x)}_{1} \\
                      &= \frac{|x_1| + |x_2|}{|x_1+x_2|}
  \end{align*}
  Daraus folgt: Die Addition zweier Zahlen mit gleichem Vorzeichen ergibt
  \begin{gather*}
    \kappa_{rel} = 1
  \end{gather*}
  Die Subtraktion zweier annähernd gleich großer  Zahlen ergibt eine sehr schlechte relative
  Konditionierung:
  \begin{gather*}
    \kappa_{rel} \gg 1
  \end{gather*}
\end{Bspe}

\textbf{Zum Beispiel} in \ref{3.1.8}: Es ist 
\begin{align*}
  x &= \begin{pmatrix}
    9,995 \\
    -9,984
  \end{pmatrix}
  \cdot 10^{-1} \\
  \widetilde{x} = fl(x) &= \begin{pmatrix}
    1 \\
    -9,98\cdot 10^{-1}
  \end{pmatrix}
  \intertext{also}
  \frac{|f(\widetilde{x})-f(x)|}{|f(x)|}	&= \frac{0,9}{1,1} 
                                                  = 0,\overline{81} \\
    &\leq \kappa_{rel}(f,x)\cdot \frac{\|\widetilde{x}-x\|_1}{\|x\|_1} \\
    &= \kappa_{rel}(f,x) \cdot 4,6\cdot 10^{-4}
\end{align*}
% 

% \subsectione{Beispiel: Lösen eines Gleichungssystems} 
\begin{Bspe}[Lösen eines lin. Gleichungssystems]
  \label{3.2.10}
  Sei $A\in \R^{n\times n}$ invertierbar und $b\in \R^n$. Es soll 
  \begin{gather*}
    Ax =b
  \end{gather*}
  gelöst werden.
  Die möglichen Lösungen in $A$ und in $b$ lassen sich folgendermaßen ermitteln:
  \begin{enumerate}[a)]
  \item Betrachte die Störungen in $b$:\\
    Sei hierzu
    \begin{gather*}
      f: b\mapsto x= A^{-1}b 
    \end{gather*}
    Berechne dann $ \kappa(f,b)$ und löse 
    \begin{align*}
      A(x + \Delta x) &= b+\Delta b \\
      f(b + \Delta b) - f(b) &= \Delta x \\
                      &= A^{-1} \cdot \Delta b && \text{da }x = A^{-1}b \\
      \Rightarrow \|\Delta x\|_{b}  &= \|A^{-1}\Delta b\|_{b} \\
                      &\leq \|A^{-1}\|_{a,b}\cdot \|\Delta b\|_{b} && \forall b, \Delta b 
    \end{align*}
    wobei $\|\cdot\| $ auf $\Renn$ die dem $\Ren$ zugeordnete Matrix-Norm sei. \\
    Die Abschätzung ist \textbf{scharf}\index{scharf}, 
    d.h. es gibt ein $\Delta b\in \R^n$, so dass \enquote{=} gilt, 
    nach Definition \ref{3.2.4}. \\
    Also gilt\footnote{vgl. auch Lemma \ref{3.2.8}: $\kappa_{abs}(f,b)=\nn{Df(b)}_{a,b}=\nn{A^{-1}}_{a,b}$}:
    \begin{gather}
      \kappa_{abs}(f,b) = \nn{A^{-1}}_{a,b} \label{III.2.6}
    \end{gather}
    unabhängig von b.  $ \quad \left( x\mapsto Ax \quad \kappa_{abs}\right)$\\
    Ebenso folgt die scharfe Abschätzung 
    \begin{align}
      \nonumber
      \frac{\|	f(b + \Delta b) - f(b)\|}{\|f(b)\|} &= \frac{\nn{\Delta x}}{\nn{x}}\\ \nonumber
                                                    & = \frac{\nn{A^{-1}\Delta b}}{\nn{x}} \\ \nonumber
                                                    & \leq  \frac{\|A^{-1}\|\cdot \|b\|}{\|x\|} \cdot \frac{\|\Delta b\|}{\|b\|} \\
      \intertext{Damit}
      \kappa_{rel} (f,b) &= \|A^{-1} \| \cdot \frac{\|b\|}{\|A^{-1}\cdot b\|} \label{III.2.7}
    \end{align}
    Da $\|b\| \leq \|A\|\cdot\|x\| = \|A\|\cdot \|A^{-1}b\|$ folgt:
    \begin{gather}
      \kappa_{rel}(f,b) \leq \|A\| \cdot \|A^{-1}\| \label{III.2.8}
    \end{gather}
    für alle (möglichen rechten Seiten) $b $.\\
    \ref{3.2.8} ist scharf in dem Sinne, dass es ein $\widehat{b}\in \R^n$ gibt 
    mit 
    \begin{gather*}
      \|\widehat{b}\| = \nn{A}\cdot \nn{\widehat{x}}
    \end{gather*}
    und somit
    \begin{gather*}
      \kappa_{rel}(f,\widehat{b}) = \nn{A}\cdot \| A^{-1}\|
    \end{gather*} %
    % 
  \item Betrachte die Störungen in $A$:\\
    Löse also 
    \begin{gather*}
      (A+\Delta A)(x+\Delta x) = b
    \end{gather*}
    Sei hierzu
    \begin{align*}
      f:A&\mapsto x= A^{-1}b \\
      \R^{n\times n}&\rightarrow \R^n
    \end{align*}
    und berechne $\kappa(f,A)$ mittels Ableitung $Df(A):\R^{n\times n} \rightarrow \R^n$:
    \begin{align*}
      C\mapsto Df(A) C&= \frac{d}{dt} \left((A+tC)^{-1} \cdot b\right) \Big\vert_{t=0} \\
                      & = \frac{d}{dt}\left((A+tC)^{-1}\right)\Big\vert_{t=0}\cdot b
    \end{align*}			
    Weiterhin gilt
    \begin{align}
      \frac{d}{dt} \left((A+tC)^{-1}\right) \Big\vert_{t=0} 
      &=-A^{-1}CA^{-1}, 
        \label{III.2.9}
    \end{align}
    da
    \begin{align*}
      0&= \frac{d}{dt}I \\
       &= \frac{d}{dt}\left( (A+tC)(A+tC)^{-1}\right)\\
       &= C(A+tC)^{-1} +(A+tC)\cdot \frac{d}{dt}(A+tC)^{-1} \\
      \Leftrightarrow \frac{d}{dt} (A+ tC)^{-1} 
       &= -(A+tC)^{-1} \cdot C(A+tC)^{-1} \, ,
    \end{align*}
    falls $(A+tC)$ invertierbar ist. Für ein genügend kleines $t$ ist das gewährleistet, da $A$ invertierbar ist (s. Lemma \ref{3.2.12}).
    \begin{gather*}
      \Rightarrow Df(A) C = -A^{-1}CA^{-1}b
    \end{gather*}
    Somit folgt
    \begin{align}
      \nonumber
      \kappa_{abs} (f,A) &= \|Df(A)\| \\ \nonumber
                         &= \sup_{\substack{
                           C\neq 0 \\ 
      C\in \R^{n\times n}											  	
      }}
      \frac{\|A^{-1}CA^{-1}b\|}{\|C\|} \\ \nonumber
                         &\leq \sup_{\substack{
                           C\neq 0 \\ 
      C \in \R^{n\times n}														  	
      }}
      \frac{\|A^{-1}\|\cdot\|C\|\cdot\|A^{-1}b\|}{\|C\|} \\ \nonumber
                         &= \|A^{-1}\| \cdot\|x\| \\ \nonumber
                         &\leq   \|A^{-1}\|^2 \cdot\|b\| \\ \nonumber
      \kappa_{rel}(f,A)  &= \frac{\|A\|}{\|f(A)\|} \cdot \|Df(A)\| \\
                         &\leq \|A\|\cdot \|A^{-1}\| \label{III.2.10}
    \end{align}
  \item betrachte Störungen in $A$ und $b$ :
    \begin{gather*}
      (A+\Delta A)(x+\Delta x) = (b+\Delta b) 
    \end{gather*}
    Für $\kappa$ müsste $\|(A,b)\|$ festgelegt werden. Dies wird jedoch nicht betrachtet. Es gilt aber folgende Abschätzung für invertierbare Matrizen $A\in \Renn $ und Störungen
    $\Delta A \in \R^{n\times n}$ mit $\|A^{-1}\|\cdot \|\Delta A\| < 1$:
    \begin{align}
      \frac{\|\Delta x\|}{\|x\|} & \leq \|A\| \cdot \|A^{-1}\|\cdot (1- \|A^{-1}\|\cdot \|\Delta A\|) 
                                   \cdot
                                   \underbrace{\left(  \frac{\|\Delta b\|}{\|b\|} +  \frac{\|\Delta A\|}{\|A\|}  \right)}_{\neq  \frac{\|(\Delta A, \Delta b)\|}{\|(A,b)\|} }
                                   \label{III.2.11}
    \end{align}
    \begin{proof} s. Übungsblatt \end{proof}
  \end{enumerate}
\end{Bspe}

% \subsectione{Definition: Kondition einer Matrix} 
\begin{Defe}
  \index{Kondition!Matrix}
  Sei $\|\cdot\|$ eine Norm auf $\R^{n\times n} $ und $A\in \R^{n\times n}$ eine reguläre Matrix.
  Die Größe
  \begin{gather*}
    \kappa_{\|\cdot\|}(A) = cond_{\|\cdot\|} \coloneqq \|A\| \cdot \|A^{-1}\|
  \end{gather*}
  heißt \textbf{Kondition der Matrix} bzgl. der Norm ${\|\cdot\|}$. \\
  Ist  ${\|\cdot\|}$ von einer Vektor-Norm ${\|\cdot\|}_p$ induziert, bezeichnet 
  $cond_p(A)$
  die $cond_{\|\cdot\|_p}(A)$. Wir schreiben $cond(A)$ für $cond_2(A)$. \\
  $cond_{\|\cdot\|}(A) $ schätzt die relative Kondition eines linearen GLS $Ax=b$ für alle möglichen 
  Störungen in $b$ oder in $A$ ab und diese Abschätzung ist scharf. \\
  
  Es stellt sich nun die Frage: \\
  \textit{Wann existiert die Inverse der gestörten invertierbaren Matrix $A$?} \\
  Hierzu werden wir die Relationen benötigen:
  \begin{align*}
    A+\Delta A &= A (I+A^{-1}\Delta A)\\
    \intertext{und mit $C \in \Renn,\, \|C\| < 1$}
    (I-C)^{-1} &= \sum_{k=0}^{\infty}C^k \\
    \|	(I-C)^{-1} \| &\leq \frac{1}{1-\|C\|}
  \end{align*}
\end{Defe}

%-----------------------------------------------------------------

% \subsectione{Lemma (Neumannsche Reihe)}
\begin{Leme}{Neumannsche Reihe}
\marginpar{27.10.2014}
  \index{Neumannsche Reihe}\label{3.2.12}
  \addtocounter{equation}{1}
  Sei $C\in\Renn$ mit $\|C\|<1$ und mit einer submultiplikativen Norm $\|\cdot\|$,
  so ist $(I-C)$ invertierbar und es gilt:
  \begin{gather*}
    (I-C)^{-1}=\sum_{k=0}^{\infty}C^k
  \end{gather*}
  Weiterhin gilt:
  \begin{gather*}
    \|(I-C)^ {-1}\| \leq \frac{1}{1-\|C\|}
  \end{gather*}
\end{Leme}

\begin{proof}
  Es gilt zu zeigen, dass $\sum_{k=1}^{\infty}C^k$ existiert: \\
  Sei $q\coloneqq \|C\| < 1$, dann gilt: 
  \begin{align*}
    \nn{ \sum_{k=0}^{m} C^k } &\leq \sum_{k=0}^{m} \nn{C^k }  && \text{Dreiecksungleichung} \\
                              &\leq \sum_{k=0}^{m}\nn{C}^k && \text{da $\nn{\,\cdot\,}$ submultiplikativ}\\
                              &=\sum_{k=0}^{m}q^k  \\
                              &= \frac{1-q^{m+1}}{1-q} \\
                              &\leq \frac{1}{1-\nn{C}} && \forall m\in \N, \text{ da } q<1 \text{ (geometr. Reihe)}
  \end{align*}
  Daraus folgt bereits, dass $\sum_{k=1}^{\infty}C^k$ existiert (nach Majorantenkriterium).\\
  Weiter gilt dann:
  \begin{align*}
    (I-C) \sum_{k=1}^{\infty}C^k &= \lim\limits_{m\rightarrow \infty}(I-C)   \sum_{k=1}^{m}C^k \\
                                 &= \lim\limits_{m\rightarrow \infty} (C^0-C^{m+1}) \\
                                 &=I 
  \end{align*}
\end{proof}

% \subsectione{Bemerkung}
\begin{Beme}
  \label{3.2.13}~
  \begin{enumerate}[a)]
  \item Für symmetrische, positiv definite Matrix $A\in \Renn$ gilt\footnote{Beweis: siehe Übungsblatt 3}: 
    \begin{gather}
      \kappa_2(A) = \frac{\lambda_{max}}{\lambda_{min}} \label{III.2.13}
    \end{gather}
  \item Eine andere Darstellung von $\kappa(A)$ ist
    \begin{gather}
      \kappa(A) \coloneqq 
      \frac{\underset{\|x\|=1}{\max}\|Ax\|}{\underset{\|x\|=1}{\min}\|Ax\|} \in  \left[ 0, \infty \right]
      \label{III.2.14}
    \end{gather}
    Diese ist auch für nicht invertierbare und rechteckige Matrizen wohldefiniert. \\
    Dann gilt offensichtlich:
  \item $\kappa(A) \geq 1$
  \item $\kappa(\alpha A)=\kappa(A) \quad \text{für } 0\neq\alpha\in\R$ (skalierungsinvariant)
  \item $A\neq 0$ und $A\in\Renn $ ist genau dann singulär, wenn $\kappa(A)=\infty$. \\
    Wegen der Skalierungsinvarianz ist die Kondition zur Überprüfung der Regularität von $A$ 
    besser geeignet als die Determinante.
  \end{enumerate}
\end{Beme}


% \subsectione{Beispiel: Kondition eines nichtlinearen Gleichungssystems}
\begin{Bspe}[Kondition eines nichtlin. Gleichungssystems]
  Sei $f:\Ren\rightarrow\Ren$ stetig differenzierbar und $y\in\Ren$ gegeben. \\
  Löse
  \begin{gather*}
    f(x) = y
  \end{gather*}
  Gesucht:
  \begin{gather*}
    \kappa(f^{-1},y)
  \end{gather*}
  mit $f^{-1}$ Ausgabe und $y$ Eingabe. \\
  Sei $Df(x)$ invertierbar, dann existiert aufgrund des Satzes für implizite Funktionen die inverse Funktion $f^{-1}$ lokal in einer Umgebung von $y$ mit $f^{-1}(y)=x$, sowie
  \begin{gather*}
    D(f^{-1})(y) = (Df(x))^{-1}
  \end{gather*}
  Hiermit folgt:
  \begin{align}
    \nonumber
    \kappa_{abs}(f^{-1},y) &= \|(Df(x))^{-1}\| \\
    \kappa_{rel}(f^{-1},y) &= \frac{\|f(x)\|}{\|x\|}\cdot\|(Df(x))^{-1}\|  \label{III.2.15}
  \end{align}
  Für skalare Funktionen $f:\R\rightarrow\R$ folgt somit:
  \begin{gather*}
    \kappa_{rel}(f^{-1},y) = \frac{|f(x)|}{|x|}\cdot \frac{1}{|f'(x)|}
  \end{gather*}
  Falls $|f'(x)|\rightarrow 0$ ist es eine schlechte absolute Kondition. \\
  Für $|f'(x)| \gg 0$ ist es eine gute absolute Kondition.\\
  \begin{figure}
    \parbox{\linewidth}{
      \centering
      \includegraphics[width=2cm]{images/image_missing.jpg}
    }
    \caption{Gute und schlechte Kondition}
  \end{figure}
  
  Damit bedeutet eine kleine Störung in $y$ eine große Störung in $x$.
\end{Bspe}


% \minisec{\Large3.2 b) Komponentenweise Konditionsanalyse} \label{3.2b}\vspace{1eM}
\extrasection{b)}{Komponentenweise Konditionsanalyse}

% \subsectione{Beispiel}
\begin{Bspe}
  Falls $A$ Diagonalgestalt hat, sind die Gleichungen unabhängig voneinander (entkoppelt)\index{entkoppelt}.
  Die erwartete relative Kondition wäre dann -- wie bei skalaren Gleicungen -- stets gleich 1.
  Ebenso sind Störungen nur in der Diagonale zu erwarten. Jedoch:
  \begin{align*}
    A  &=\begin{pmatrix}
      1 & 0\\
      0 & \varepsilon
    \end{pmatrix} \\
    \Rightarrow 	A^{-1}&=\begin{pmatrix}
      1 & 0\\
      0 & \varepsilon^{-1}
    \end{pmatrix}\\
    \Rightarrow \kappa_\infty& = \kappa_2 = \frac{1}{\varepsilon} 
        && \text{für }0 < \varepsilon \leq 1											
  \end{align*}
\end{Bspe}

% \subsectione{Definition: Komponentenweise Kondition}
\begin{Defe}
  \index{Kondition!komponentenweise}
  Sei $(f, x) $ ein Problem mit $f(x)\neq 0$ und $x=(x_i)_{i=1,\cdots , n}$ mit $x_i\neq 0 $  für alle $i=1,\cdots, n$.
  Die \textbf{komponentenweise Kondition} von $(f,x) $ ist die kleinste Zahl $\kappa_{rel}\geq 0$, so dass:
  \begin{align*}
    \frac{\|f(\widetilde{x})-f(x)\|_\infty}{\|f(x)\|_\infty} 
    &\leq \kappa_{rel} \cdot \underset{i}{\max}\frac{|\widetilde{x_i}-x_i|}{|x_i|}+ o\left(\underset{i}{\max}\frac{|\widetilde{x_i}-x_i|}{|x_i|}\right) 
    && \text{für }\widetilde{x}\rightarrow x
  \end{align*}
  Vorsicht:
  \begin{gather*}
    \frac{\|\widetilde{x}-x\|_\infty}{\|x\|_\infty}\neq \underset{i}{\max}\frac{|\widetilde{x_i}-x_i|}{|x_i|}
  \end{gather*}
\end{Defe}

% \subsectione{Lemma} 
\begin{Leme}
  \label{3.2.17}
  Sei $f$ differenzierbar und fasse $|\cdot|$ komponentenweise auf, d.h. $|x| = \begin{pmatrix}
    |x_1| \\
    \vdots \\
    |x_n|
  \end{pmatrix}$.
  Dann gilt:
  \begin{gather}
    \kappa_{rel} = \frac{\|\, |Df(x)|\cdot |x| \, \|_\infty}{\|f(x)\|_\infty} \label{III.2.16}
  \end{gather}
\end{Leme}

\begin{proof}
  Vergleiche seien ebenfalls komponentenweise zu verstehen. \\
  Nach dem Satz von Taylor gilt: 
  \begin{align*}
    f_i(\widetilde{x})-f_i(x) 
    &= \left( \frac{\partial f_i}{\partial x_i}(x), \cdots ,\frac{\partial f_i}{\partial x_n}(x) \right)
      \cdot \begin{pmatrix}
        \widetilde{x}_1-x_1 \\
        \vdots \\
        \widetilde{x}_n-x_n
      \end{pmatrix}
    + o\left(\|\widetilde{x}-x\|\right) \\
    \Rightarrow |f_i(\widetilde{x})-f_i(x)|
    &\leq |Df(x)|
      \cdot \begin{pmatrix}
        |x_1|\cdot \frac{\widetilde{x}_1-x_1 }{|x_1|}\\
        \vdots \\
        |x_n|\cdot \frac{\widetilde{x}_n-x_n }{|x_n|}
      \end{pmatrix}
    + o\left(\underset{i}{\max}\frac{\widetilde{x}_i-x_i }{|x_i|}\right) 
    && \text{da $x_i$ fest und $\widetilde{x}_i\rightarrow x_i$} \\
    &\leq |Df(x)| \cdot |x| \cdot \underset{i}{\max}\frac{\widetilde{x}_i-x_i }{|x_i|}
      +o\left(\underset{i}{\max}\frac{\widetilde{x}_i-x_i }{|x_i|}\right) \\
    \Rightarrow \frac{\|f(\widetilde{x})-f(x)\|_\infty}{\|f(x)\|_\infty}
    &\leq  \frac{\|\, \|Df(x)|\cdot |x| \, \|_\infty}{\|f(x)\|_\infty}
      \cdot \underset{i}{\max}\frac{\widetilde{x}_i-x_i }{|x_i|}
      + o\left( \underset{i}{\max}\frac{\widetilde{x}_i-x_i }{|x_i|} \right)
  \end{align*}
  Wähle $\widetilde{x}_i = x_j+h\cdot sign \frac{\partial f_i}{\partial x_j}(x)$ mit $h>0$,
  dann gilt:
  \begin{gather*}
    |Df_i(x)(\widetilde{x}-x)| = Df_i(x)(\widetilde{x}-x)
  \end{gather*}
  und in obiger Rechnung gilt Gleichheit. \\
  Also folgt, dass
  \begin{align*}
    \frac{\|\,|Df(x)|\cdot |x| \, \|_\infty}{\|f(x)\|_\infty} &= \kappa_{rel} 
  \end{align*}
\end{proof}

% \subsectione{Beispiel}
\begin{Bspe}~
  \begin{enumerate}[a)]
  \item Komponentenweise Kondition der Multiplikation
    \begin{align*}
      f:&\R^2 \rightarrow \R, \, f(x,y) \coloneqq x\cdot y \\
      \Rightarrow Df(x,y) &= (y, x)  \\
      \Rightarrow \kappa_{rel}(x,y) &= \frac{\left\| (|y|, |x|)\cdot \begin{pmatrix}
            |x| \\
            |y|
          \end{pmatrix}\right\|_\infty}
      {|x\cdot y|} \\
        &= \frac{2\cdot|x|\cdot |y|}{|x\cdot y|} \\
        &= 2
    \end{align*}
  \item Komponentenweise Kondition eines linearen Gleichungssystems:\\
    Löse $Ax=b$ mit möglichen Störungen in $b$, also zu
    \begin{align*}
      f: & b\mapsto A^{-1}b \\
      \kappa_{rel} & = \frac{\| \, |A^{-1}| \cdot |b|\, \|_\infty}{\|A^{-1}b\|_\infty}
    \end{align*}
    Falls A eine Diagonalmatrix ist, folgt:
    \begin{gather*}
      \kappa_{rel}=1
    \end{gather*}
  \item Komponentenweise Kondition des Skalarproduktes:
    \begin{align*}
      \langle x,y \rangle \coloneqq \sum_{i=1}^{n}x_i y_i& = x^Ty \\
      f: \R^2 \rightarrow \R, \, f(x,y) &= \langle x,y \rangle \\
      \Rightarrow Df(x,y) &= (y^T, x^T) \\
      \kappa_{rel}  &= \frac{\left\| \,\left|(y^T, x^T)\right|\cdot\left|\begin{pmatrix}
              x \\
              y
            \end{pmatrix}\right|\right\|_\infty }
      {\|\langle x,y\rangle\|_\infty}\\
                                                         &= \frac{2\cdot |y^T|\cdot |x|}{|\langle x,y\rangle|} \\
                                                         &= 2\cdot \frac{\langle |x|,|y|\rangle}{|\langle x,y\rangle|} \\
                                                         &= 2 \cdot \frac{\cos(|x|, |y|)}{\cos(x,y)}  \\
                                                         &&&				\text{	da  }\cos(x,y) = \frac{\langle y,x \rangle}{\|x\|_2 \cdot \|y\|_2} \, . 
    \end{align*}
    Falls $x$ und $y$ nahezu senkrecht aufeinander stehen, kann das Skalarprodukt sehr schlecht konditioniert sein. \\
    Zum Beispiel für $x=\widetilde{x} = \begin{pmatrix} 1 \\1 \end{pmatrix}$
    und $y=\begin{pmatrix} 1+10^{-10} \\-1 \end{pmatrix},
    \, \widetilde{y}=\begin{pmatrix} 1 \\-1 \end{pmatrix}$. \\
    \begin{figure}
      \parbox{\linewidth}{
        \centering
        \includegraphics[width=2cm]{images/image_missing.jpg}
      }
      \caption{Schlechte Kondition des Skalarprodukts bei nahezu senkrechten Vektoren}
    \end{figure}
  \end{enumerate}	
\end{Bspe}

\sectione{Stabilität von Algorithmen}
Bislang: Kondition eines gegebenen Problems $(f,x)$. \\
Nun stellt sich die Frage: \textit{Was passiert durch das Implementieren am Rechner? }\\
Ein \enquote{stabiler} Algorithmus sollte ein gut konditioniertes Problem nicht \enquote{kaputt machen}.\\

\begin{figure}
  \parbox{\linewidth}{
    \centering
    \includegraphics[width=2cm]{images/image_missing.jpg}
  }
  \caption{Stabilität eines Algorithmus}
\end{figure}

% \minisec{\Large3.3 a) Vorwärtsanalyse} \label{3.3a}\vspace{1eM}
\extrasection{a)}{Vorwärtsanalyse}

Die Fehlerfortpflanzung durch die einzelnen Rechenschritte, aus denen die Implementierung aufgebaut ist, wird abgeschätzt.

% \subsectione{Bemerkung}
\begin{Beme}
  Für die Rechenoperationene $+,-,\, \cdot \, , \, /\,$, kurz $\nabla$, gilt:
  \begin{align}
    \nonumber
    fl(a\nabla b) &= (a\nabla b)\cdot (1+\varepsilon) \\
                  &= (a\nabla b) \cdot \frac{1}{1+\mu} \label{III.3.1}
  \end{align}
  mit $|\varepsilon|, |\mu| \leq eps$.
\end{Beme}

%--------------------------------------------------------------

% \subsectione{Beispiel}
\begin{Bspe}
\marginpar{29.10.2014}
  Sei $f(x_1, x_2, x_3) \coloneqq \frac{x_1x_2}{x_3}$ mit Maschinenzahlen $x_i$ und $x_3\neq 0$ und sei der Algorithmus durch
  \begin{gather*}
    f(x_1, x_2, x_3) = (f^{(2)} \circ f^{(1)})(x_1, x_2, x_3) 
  \end{gather*}
  gegeben mit 
  \begin{align*}
    f^{(1)}(x_1, x_2, x_3) & = (x_1\cdot x_2, x_3) && \text{und} \\
    f^{(2)}(y,z) &= \frac{y}{z}
  \end{align*}
  Die Implementierung $\widetilde{f}$ von $f$  beinhaltet Rundungsfehler. \\

  Sei  $x=(x_1, x_2, x_3) $. Daraus folgt:
  \begin{align*}
    \widetilde{f}^{(1)}(x) &= (fl(x_1\cdot x_2), x_3) \\
                           & = (x_1x_2 (1+\varepsilon_1), x_3)
                             \intertext{mit $|\varepsilon_1|\leq eps$:}
                             \widetilde{f}(x) &= \widetilde{f}^{(2)}(\widetilde{f}^{(1)}(x)) \\
                           &= fl(f^{(2)}(x_1 x_2 (1+\varepsilon_1), x_3)) \\
                           &= \frac{x_1x_2(1+\varepsilon_1)}{x_3}\cdot (1+\varepsilon_2)  \\
                           &= f(x)\cdot (1+\varepsilon_1)(1+\varepsilon_2)
                             \intertext{mit $|\varepsilon_2| \leq eps$:}
                             \frac{|\widetilde{f}(x) -f(x)|}{|f(x)|} &= |\varepsilon_1+\varepsilon_2 +\varepsilon_1\cdot \varepsilon_2| \\
                           &\leq 2eps + eps^2
  \end{align*}
  Dies ist eine \enquote{worst case} Analyse, da immer der maximale Fehler angenommen wird,
  und gibt i.d.R. eine starte Überschätzung des Fehlers an.
  Für qualitative Aussagen sind sie jedoch unnützlich. \\
  In Computersystemen stehen mehr Operationen wie $\nabla$ zur Verfügung,
  die mit einer relativen Genauigkeit $eps$ realisiert werden können.	
\end{Bspe}

Daher:

% \subsectione{Definition: Elementar ausführbar}
\begin{Defe}
  Eine Abbildung $\phi : U\subseteq \Ren \rightarrow \R^m$ heißt
  \textbf{elementar ausführbar}\index{elementar ausführbar}, falls es 
  eine elementare Operation $\widetilde{\phi}:\F^n \rightarrow \F^m$
  gibt, wobei $\F$ die Menge der Maschinenzahlen bezeichne mit
  \begin{gather}
    |\widetilde{\phi}_i(x)-\phi_i(x)| \leq eps\cdot |\phi_i(x) | 
    \quad \forall x\in \F^n \text{ und } i=1,\cdots , m \label{III.3.2}\, .
  \end{gather}
  $\widetilde{\phi}$ heißt dann \textbf{Realisierung}\index{Realisierung} von $\phi$.
\end{Defe}


\textbf{Bemerkung}\\
aus \eqref{III.3.2} folgt für $1\leq p\leq \infty$:
\begin{gather}
  \nn{\widetilde{\phi}(x)-\phi(x)}_p \leq eps\cdot\nn{\phi(x)}_n 
  \quad \forall x\in\F^n \label{III.3.3}
\end{gather}


% \subsectione{Definition: Algorithmus, Implementation}
\begin{Defe}
  Sei $f:E\subseteq \Ren \rightarrow \R^m$ gegeben.\\
  Ein Tupel $\left(f^{(1)},\cdots ,f^{(l)}\right)$ mit $l\in \N$ von elementar ausführbaren
  Abbildungen
  \begin{gather*}
    f^{(i)}: U_1\subseteq \R^{k_i} \rightarrow U_{i+1}\subseteq \R^{k_{i+1}}
  \end{gather*}
  mit $k_1=n$ und $k_{l+1}=m$ heißt \textbf{Algorithmus}\index{Algorithmus} von $f$, falls
  \begin{gather*}
    f=f	^{(l)}\circ \dotsc \circ f^{(1)}
  \end{gather*}
  Das Tupel $(\widetilde{f}1^{(1)},\cdots ,\widetilde{f}^{(l)})$ mit Abbildungen $\widetilde{f}^{(i)}$, welche Realisierungen der $f^{(i)}$ sind,
  heißt \textbf{Implementation}\index{Implementation} von 
  $\left(f^{(1)},\dotsc ,f^{(l)}\right)$.
  Die Komposition 
  \begin{gather*}
    \widetilde{f}=\widetilde{f}	^{(l)}\circ \dotsc \circ \widetilde{f}^{(1)}
  \end{gather*}
  heißt Implementation von f. \\
  Im Allgemeinen gibt es verschiedene Implementierungen einer Abbildung $f$.
\end{Defe}


% \subsectione{Lemma (Fehlerfortpflanzung)}
\begin{Leme}[Fehlerfortpflanzung]
  \label{3.3.5} \index{Fehler!Fortpflanzung}
  Sei $x\in \Ren$ und $\widetilde{x}\in \F^n$ mit $|\widetilde{x}_i-x_i|\leq eps|x_i|$ für alle 
  $i=1,\cdots , n$.
  Sei $\left(f^{(1)},\dotsc ,f^{(l)}\right)$ ein Algorithmus für $f$ und 
  $(\widetilde{f}^{(1)},\dotsc ,\widetilde{f}^{(l)})$ eine zugehörige Implementation. \\
  Mit den Abkürzungen
  \begin{align*}
    x^{(j+1)} &\coloneqq f^{(j)}\circ \dotsc \circ f^{(1)}(x) \\
    x^{(1)} &\coloneqq x
  \end{align*}
  und entsprechend mit $\widetilde{x}^{(j+1)}$ gilt,
  falls $x^{(j+1)} \neq 0$ für alle $j=0,\dotsc , (l-1)$ und $\nn{\,\cdot\,}$ eine beliebige p-Norm ist:
  \begin{align}
    \frac{\nn{\widetilde{x}^{(j+1)}-x^{(j+1)}}}{\nn{x^{(j+1)}}}
    &\leq eps \cdot \K + o\left(eps\right)
      \label{III.3.4} 
    \\ \nonumber
    \K^{(j)}&=(1+\kappa^{(j)}+\kappa^{(j)}\cdot \kappa^{(j-1)}+ \cdots + \kappa^{(j)}\cdot \dotsm \cdot \kappa^{(1)}) \\ \nonumber
  \end{align}
  wobei $	\kappa^{(j)} \coloneqq \kappa_{rel}(f^{(j)}, x^{(j)})$ die Kondition der elementar ausführbaren Operationen $f^{(j)}$ ist.
\end{Leme}

\begin{proof}
  \begin{align*}
    \frac{\nn{\widetilde{x}^{(j+1)}-x^{(j+1)}}}{\nn{x^{(j+1)}}}
    &= \frac{\nn{\widetilde{f}^{(j)}(\widetilde{x}^{(j)})-f^{(j)}(x^{(j)})}}
      {\nn{f^{(j)}(x^{(j)})}} \\
    &\leq \frac{\nn{\widetilde{f}(\widetilde{x})-f(\widetilde{x})}}{\nn{f(\widetilde{x})}}
      \cdot \frac{\nn{f(\widetilde{x})}}{\nn{f(x)}}
      + \frac{\nn{f(\widetilde{x})-f({x})}}{\nn{f(x)}} 
      \quad\quad \text{(Index $j$ vernachlässigt)}\\
    &\leq eps \left( 1+ \frac{\nn{f(\widetilde{x})-f({x})}}{\nn{f(x)}}\right)
      + \frac{\nn{f(\widetilde{x})-f({x}})}{\nn{f(x)}}\\
    &\overset{\text{nach \ref{III.3.3}}}{=} eps + (eps+1) \cdot 						\left(\kappa{(j)}\cdot \frac{\nn{\widetilde{x}^{(j)}-x^{(j)}}}{\nn{x^{(j)}}}\right)
      + o\left( \frac{\nn{\widetilde{x}^{(j)}-x^{(j)}}}{\nn{x^{(j)}}}\right)
  \end{align*}
  Nach Voraussetzung gilt Gleichung \eqref{III.3.4}  mit $\K^{(0)}=1$ für $j=0$. \\
  Für $j=1$ folgt nach Voraussetzung mit Gleichung \eqref{III.3.3}
  \begin{align*}
    \frac{\nn{\widetilde{x}^{(2)}-x^{(2)}}}{\nn{x^{(2)}}}
    & \leq eps +(eps+1) \cdot \left( \kappa^{(1)}eps+ o(eps)\right) \\
    &= eps(1+\kappa^{(1)}) + o(eps) \\
    &= eps\K^{(1)} + o(eps)
  \end{align*}
  Womit der Induktionsanfang gezeigt ist. \\
  Für den Induktionsschritt von $j-1$ zu $j$:
  \begin{align*}
    \frac{\nn{\widetilde{x}^{(j+1)}-x^{(j+1)}}}{x^{(j+1)}}
    & \leq eps + (1+eps)\kappa^{(j)} \left[ eps \K^{(j-1)}+ o(eps) \right] \\
    &\phantom{\leq eps+} + (1+eps) \cdot o\left( eps\cdot \K^{(j-1)} +o(eps)\right) \\
    &= eps\left(1+\kappa^{(j)}\cdot \K^{(j-1)}\right)+ o(eps)
  \end{align*}
  Mit $\K^{(j)} = 1+ \kappa^{(j)}\cdot \K^{(j-1)}$ folgt die Behauptung.
\end{proof}

Hiermit folgt:

% \subsectione{Korollar}
\begin{Kore}
  \label{3.3.6}
  Unter der Voraussetzung von Lemma \ref{3.3.5} gilt:
  \begin{gather}
    \frac{\nn{\widetilde{f}(\widetilde{x})-f(x)}}{\nn{f(x)}} \leq 
    eps\cdot \left( 1+\kappa^{(l)}+ \kappa^{(l)}\cdot \kappa^{(l-1)}+ \dotsc
      + \kappa^{(l)}\cdot \dotsc \cdot \kappa^{(1)}\right) + o(eps) 
    \label{III.3.5}
  \end{gather}~
\end{Kore}

% \subsectione{Bemerkung}
\begin{Beme}
  Mit Korollar \ref{3.3.6} ist offensichtlich, dass schlecht konditionierte Probleme 
  zu elementar ausführbaren Abbildungen so früh wie möglich ausgeführt werden sollten. \\
  Nach Beispiel \ref{3.2.9} ist die Substraktion zweier annähernd gleicher Zahlen schlecht konditioniert.
  Deshalb sollte man unvermeidbare Subtraktionen möglichst früh durchführen. \\
  Allerdings hängt $\kappa^{(j)}$ nicht nur von $f^{(j)}$, sondern auch vom 
  Zwischenergebnis $x^{(j)}$ ab, welches a priori unbekannt ist.
\end{Beme}

% \subsectione{Bemerkung zur Sprechweise} %\label{3.3.8}
\begin{Beme}[Sprechweise]
  Der Quotient 
  \begin{gather}
    \frac{\overbrace{\frac{\nn{\widetilde{f}(\widetilde{x})-f(x)}}{\nn{f(x)}}}^{
        \text{Gesamtfehler}}}
    {\underbrace{\frac{\nn{\widetilde{f}(\widetilde{x})}}{\nn{f(x)}}}_{
        \scriptsize\substack{
          \text{Fehler} \\
          \text{durch} \\
          \text{Problem}
        }}
      \cdot
      {\underbrace{\frac{\nn{\widetilde{x}-x}}{\nn{x}}}_{
          \scriptsize\substack{\text{Eingabe-} \\ \text{fehler}}}}}
    \label{III.3.6}
  \end{gather}
  gibt die \textbf{Güte des Algorithmus} \index{Güte!Algorithmus} an.
  Als Stabilitätsindikator kann also 
  \begin{gather}
    \sigma\left(f, \widetilde{f}, x\right) \coloneqq \frac{\K}{\kappa_{rel}(f, x)}
    \label{III.3.7}
  \end{gather}
  verwendet werden und es gilt
  \begin{gather*}
    \frac{\nn{\widetilde{f}(\widetilde{x})-f(x)}}{\nn{f(x)}}
    < \underbrace{\sigma\left( f,\widetilde{f}, x\right) }_{
      \substack{\text{Beitrag}\\
        \text{des} \\
        \text{Algorithmus}}}
    \cdot \underbrace{\kappa_{rel}(f,x)}_{
      \substack{\text{Beitrag} \\
        \text{des} \\
        \text{Problems}}}
    \cdot \underbrace{eps}_{\substack{\text{Rundungs-}\\\text{fehler}}}
    + \quad o(eps)
  \end{gather*}
  Falls $\sigma( f,\widetilde{f}, x)  < 1$, dämpft der Algorithmus die Fehlerfortpflanzung der Eingabe- und Rundungsfehler und heißt \textbf{stabil}\index{Stabilität}. \\
  Für $\sigma( f,\widetilde{f}, x)  \gg 1$ heißt der Algorithmus \textbf{instabil}.
\end{Beme}



% \subsectione{Beispiel}
\begin{Bspe}
  Nach Gleichung \eqref{III.3.3} gilt für die Elementaroperationen $\K\leq 1$.
  Da für die Subtraktion zweier annähernd gleich großer Zahlen $\kappa_{rel}\gg 1$ gilt,
  ist der Stabilitätsfaktor zweier annähernd gleich großer
  Zahlen sehr klein und der Algorithmus also stabil,
  Falls es sich jedoch bei einer zusammengesetzten Abbildung $f=h\circ g$
  bei der zweiten Abbildung $h$ um eine Subtaktion handelt, gilt
  \begin{gather*}
    \K =(1+\kappa(sub)+\kappa(sub)\cdot\kappa(g))
  \end{gather*}
  und die Stabilität ist gefährdet.
  Genauere Abschätzungen und damit genauere Indikatoren
  können durch komponentenweise Betrachtungen erhalten werden.
\end{Bspe}


\extrasection{b)}{Rückwärtsanalyse} \index{Rückwärtsanalyse}
Die Fragestellung ist nun: \\
\textit{Kann $\widetilde{f}(\widehat{x})$ als exaktes Ergebnis von einer gestörten Eingabe $\widehat{x}$ unter der exakten Abbildung $f$ aufgefasst werden?}\\
Das würde heißen
\begin{gather*}
  \exists\, \widehat{x}\in \Ren: f(\widehat{x})= \widetilde{f}(\widetilde{x}) \, .
\end{gather*}
Dann schätze den Fehler $ \nn{\widehat{x}-x}$
bzw. für nicht injektive $f$
\begin{gather*}
  \min_{\widehat{x}\in \Ren}
  \left\{
    \nn{\widehat{x}-x} 
    \middle\vert f(\widehat{x}) = \widetilde{f}(\widetilde{x}) 
  \right\}
\end{gather*} 
ab. \\

\begin{tikzpicture}[scale=1,line cap=round,line join=round,>=triangle 45,x=1.0cm,y=1.0cm]
  \matrix (m) [matrix of math nodes,row sep=0.4cm,column sep=3cm,minimum width=3em]
  {
    A &L,R & \parbox{3cm}{$x=A^{-1}b$}\\
    \hat{A} &\hat{L},\hat{R} &\parbox{3cm}{$ \hat{x}=\left(\hat{L}\hat{R}\right)^{-1}b$}\\
    \overline{A}&&\parbox{3cm}{$\overline{x}=\left(\overline{A}\right)^{-1}b$}\\
  };
  \path
  (m-1-1) edge [->](m-1-2)
  edge [->,decorate,decoration=zigzag](m-2-2)
  (m-1-2) edge [->](m-1-3)
  (m-2-1) edge [<-] (m-2-2)
  (m-2-2) edge [->] (m-2-3)
  edge [->,decorate,decoration=zigzag](m-3-3)
  (m-3-1) edge [<-] (m-3-3);
\end{tikzpicture}

Ein Anwendungsbeispiel:\\
Die Eingangsdaten seien Messdaten $\tilde{x}$  mit 1 \% relativer Genauigkeit. 
Liefert die Rückwärtsanalyse, dass $\tilde{f} (\tilde{x})$ als exaktes Ergebnis 
$f(\hat x)$ mit Eingangsdaten $\hat x$, 
die höchstens um 0,5 \% schwanken, aufgefasst werden kann,
so ist das Verfahren \enquote{geeignet} . 

Die Rückwärtsanalyse ist
\begin{itemize}
\item in der Regel leichter durchführbar als die
  Vorwärtsanalyse  und 
\item ebenfalls nur eine qualitative Schätzung der
  Genauigkeit der numerisch berechneten Werte.
\end{itemize}

% \subsectione{Bemerkung}
\begin{Beme}
  $$ \mbox{Vorw"artsfehler} \; \leq \; \mbox{Kondition des
    Problems } \cdot \mbox{R"uckw"artsfehler} .$$
  
  $$ ||\tilde f (\tilde x)  - f(x)|| \leq \kappa (f,x) \; || \hat x -x|| $$
  
  Beispiel: Rückwärtsanalyse der Gauß-Elimination
  (geht auf Wilkinson zurück)
\end{Beme}


% \subsectione{Satz}
\begin{Satze}
  $A \in \R^{n\times n}$ besitze eine LR-Zerlegung. Dann berechnet die
  Gauß-Elimination Matrizen $\hat{L}$ und $\hat{R}$,
  so dass
  \begin{gather*}\hat{L}\,\hat{R}\; = \;\hat{A}\end{gather*}
  und
  \begin{align*}
    \nn{\hat{ A}-A} & \leq 
    \frac{eps}{1-n \cdot eps}
      \left( \nn{\hat{L}}
      \nn{\begin{pmatrix}
               1 && & 0\\
               &2\\
               &&\ddots \\
                0&&& n\\
       \end{pmatrix}}\nn{\hat{R}}-\nn{\hat{R}}
       \right)\\
      & \leq  \frac{n\cdot eps}{1-n\cdot eps}\nn{\hat{L}}\nn{\hat{R}} \\
      & = n\cdot eps \nn{\hat{L}}\cdot \nn{\hat{R}\,} +
                                 \mathcal{O}( n^2 eps^2)
  \end{align*}
  falls $n\cdot eps \leq \frac{1}{2}$.
\end{Satze}
\begin{proof}
  \cite[siehe][]{stoerbulirsch}.
\end{proof}


% \subsectione{Satz (Sautter 1971)}
\begin{Satze}[Sautter 1971]
  $ A \in \R^{n\times n}$ besitze eine LR-Zerlegung.  Dann berechnet das
  Gaußsche Eliminationsverfahren für das Gleichungssystem
  $\;  A  x =  b$ eine Lösung $\overline{ x}$ mit 
  \begin{gather*} \overline{ A} \overline{ x}  =   b \vspace*{-2ex} \end{gather*}
  mit
  \begin{gather*} 
    | \overline{ A} - A |  \leq  2 \, n \, eps \,
    |\hat{L}| \, | \hat{R}| + \mathcal{O}( n^2 eps^2).
  \end{gather*}
\end{Satze}

\begin{proof}
  \cite[siehe][]{deuflhardhohmann}.
\end{proof}

Weitere Abschätzungen existieren für Gauß-Elimination mit 
Pivotisierung und für spezielle Klassen von Matrizen.

\subsectione{Allgemeine Faustregeln für die LR-Zerlegung}
\label{III.3.13}
\begin{itemize}
\item Falls die Matrix $n | \hat{ L} | \,  |\hat{ R}|$ die
  selbe Größenordnung wie $| A|$ besitzt, ist der
  Algorithmus \enquote{gutartig};
\item Für tridiagonale Matrizen  ist der Algorithmus mit
  Spaltenpivotisierung stabil.
\item Falls $ A$ oder $ A^T$  \textbf{strikt diagonal
    dominant}\index{strikt diagonal dominant} ist, d.h. 
  \begin{gather*}
    | a_{ii} | > \sum\limits_{j=1 ,\, j \not = i}^{n} | a_{ij}| 
    \mbox{ für alle } i = 1, \ldots, n,
  \end{gather*}
  ist Spaltenpivotisierung überflüssig. Der Algorithmus ist stabil.
\item Für symmetrische, positiv definite Matrizen sollte keine Pivotisierung
  durchgeführt werden, um die Symmetrie zu erhalten.
  Der Algorithmus ist stabil.
\end{itemize}

\subsubsection{Vorsicht}
Selbst wenn die LR-Zerlegung stabil ist, in dem Sinne dass
$| \overline{ A} - A|$ klein ist für
$ \overline{ A}\overline{ x}  =  b $, 
kann die numerische Lösung $\overline{ x}$ sehr
ungenau sein, da der Vorwärtsfehler $| \overline{ x} - {  x}| $ auch von der
Kondition abhängt.

Ein Beispiel hierzu ist die Hilbertmatrix
\begin{gather*}
  H  =  \left( \frac{1}{i + j -1 } \right)_{i,j= 1,\ldots, n}\, ,
\end{gather*}
für die $cond ( H)$ exponentiell mit der Dimension $n$ wächst.


\sectione{Beurteilung von Näherungslösungen linearer GLS}
Zu $Ax=b$ liege eine Näherungslösung $\widetilde{x}$ vor.


\extrasection{a)}{Im Sinne der Vorwärtsanalyse}
Im Sinne der Vorwärtsanalyse und der Fehlerentwicklung durch das Problem gilt:
\begin{gather*}
  \frac{\nn{\widetilde{x}-x}}{\nn{x}} \leq cond(A) \cdot \frac{\nn{\Delta b}}{\nn{b}}
\end{gather*}
nach Beispiel \ref{3.2.10}, 
mit dem Residuum 
\begin{align}
  r(\widetilde{x})  & \coloneqq A\widetilde{x} - b \label{III.4.1} \\ \nonumber 
                    &	= \widetilde{b}-b \\ \nonumber
                    & = \Delta b
\end{align}
Wie der absolute Fehler ist das Residuum skalierungsabhängig.
Daher ist $\nn{r(\widetilde{x})}$ \enquote{klein} ungeeignet, um
Genauigkeitsaussagen zu treffen. \\
Um den Fehler in $x$ abzuschätzen, ist die Betrachtung von 
\begin{gather}
  \frac{\nn{r(\widetilde{x})}}{\nn{b}} \label{III.4.2}
\end{gather}
geeigneter. \\
Für große $cond(A)$ ist dieser Quotient jedoch weiterhin ungeeignet.

\extrasection{b)}{Im Sinne der Rückwärtsanalyse}
% \subsectione{Satz (Prager und Oettli, 1964)} 
\begin{Satze}[Prager und Oettli, 1964]\label{3.4.1}
  Sei $\tilde{ x}$ eine Näherungslösung für 
  $  A  x =  b$. 
  Falls
  \begin{gather}\label{III.4.3}
    |  r(\tilde { x})|  \; \leq \; \varepsilon ( | A| | \tilde { x} | + |  b|).
  \end{gather}
  dann existiert eine Matrix $\tilde{ A}$  und ein
  Vektor $\tilde { b}$, so dass
  \begin{gather*}
    \tilde{ A} \tilde { x}  =  \tilde{ b} 
  \end{gather*}
  und
  \begin{equation}
    |\tilde{ A} -  A |  \leq  \varepsilon | A|
    \quad \textrm{und} \quad | \tilde{ b} -  b| \leq
    \varepsilon | b|.
    \label{III.4.4}
  \end{equation}
  
  Aufgrund von \eqref{III.4.3} wird der komponentenweise relative
  Rückwärtsfehler durch 
  \begin{gather*}
    \max_i \frac{|  A \tilde{ x} -  b|_i}
    { (| A|\, |\tilde{ x}| + | b|)_i} 
  \end{gather*}
  abgeschätzt.
  
  Für den normweisen relativen Rückwärtsfehler gilt entsprechend
  (Rigal und Gaches 1967)
  \begin{gather*}
    \frac{\|  A \tilde{ x} -  b\|}
    { \| A\| \|\tilde{ x} \| + \| b\| } .
  \end{gather*}
\end{Satze}
